\documentclass[11pt,reqno]{amsart}
\usepackage[localtheorems]{raeez-math-template}

% ================================================================
% Packages
% ================================================================
\usepackage{array}

% ================================================================
% Theorem environments
% ================================================================
\newtheorem{theorem}{Theorem}[section]
\newtheorem{proposition}[theorem]{Proposition}
\newtheorem{lemma}[theorem]{Lemma}
\newtheorem{corollary}[theorem]{Corollary}
\theoremstyle{definition}
\newtheorem{definition}[theorem]{Definition}
\newtheorem{construction}[theorem]{Construction}
\newtheorem{example}[theorem]{Example}
\theoremstyle{remark}
\newtheorem{remark}[theorem]{Remark}

% ================================================================
% Macros
% ================================================================
\providecommand{\cA}{\mathcal{A}}
\providecommand{\cB}{\mathcal{B}}
\providecommand{\cC}{\mathcal{C}}
\providecommand{\cD}{\mathcal{D}}
\providecommand{\cE}{\mathcal{E}}
\providecommand{\cF}{\mathcal{F}}
\providecommand{\cG}{\mathcal{G}}
\providecommand{\cH}{\mathcal{H}}
\providecommand{\cK}{\mathcal{K}}
\providecommand{\cL}{\mathcal{L}}
\providecommand{\cM}{\mathcal{M}}
\providecommand{\cN}{\mathcal{N}}
\providecommand{\cO}{\mathcal{O}}
\providecommand{\cP}{\mathcal{P}}
\providecommand{\cR}{\mathcal{R}}
\providecommand{\cS}{\mathcal{S}}
\providecommand{\cW}{\mathcal{W}}
\providecommand{\cZ}{\mathcal{Z}}
\providecommand{\barB}{\bar{B}}
\providecommand{\fg}{\mathfrak{g}}
\providecommand{\fh}{\mathfrak{h}}
\providecommand{\fsl}{\mathfrak{sl}}
\providecommand{\bC}{\mathbb{C}}
\providecommand{\bQ}{\mathbb{Q}}
\providecommand{\bR}{\mathbb{R}}
\providecommand{\bZ}{\mathbb{Z}}
\providecommand{\bN}{\mathbb{N}}
\providecommand{\CC}{\mathbb{C}}
\providecommand{\ZZ}{\mathbb{Z}}
\providecommand{\RR}{\mathbb{R}}
\providecommand{\QQ}{\mathbb{Q}}
\providecommand{\HHH}{\mathbb{H}}
\providecommand{\Ran}{\mathrm{Ran}}
\providecommand{\MC}{\mathrm{MC}}
\providecommand{\Sym}{\mathrm{Sym}}
\providecommand{\Hom}{\mathrm{Hom}}
\providecommand{\Ext}{\mathrm{Ext}}
\providecommand{\Res}{\mathrm{Res}}
\providecommand{\Aut}{\mathrm{Aut}}
\providecommand{\FM}{\overline{C}}
\providecommand{\Ass}{\mathrm{Ass}}
\providecommand{\Com}{\mathrm{Com}}
\providecommand{\Vir}{\mathrm{Vir}}
\providecommand{\Conf}{\mathrm{Conf}}
\providecommand{\FP}{\mathrm{FP}}
\providecommand{\conn}{\mathrm{conn}}
\providecommand{\ord}{\mathrm{ord}}
\providecommand{\sh}{\mathrm{sh}}
\providecommand{\av}{\mathrm{av}}
\providecommand{\SK}{\mathrm{SK}}
\providecommand{\Sp}{\mathrm{Sp}}
\providecommand{\Id}{\mathrm{Id}}
\providecommand{\id}{\mathrm{id}}
\providecommand{\gmod}{\mathfrak{g}^{\mathrm{mod}}_\cA}
\providecommand{\Defcyc}{\mathrm{Def}_{\mathrm{cyc}}}
\providecommand{\ChirHoch}{\mathrm{ChirHoch}}

\DeclareMathOperator{\gr}{gr}
\DeclareMathOperator{\rk}{rk}
\DeclareMathOperator{\Spec}{Spec}
\DeclareMathOperator{\tr}{tr}

\numberwithin{equation}{section}

% ================================================================
\begin{document}

\title[Arithmetic Shadows of Chiral Algebras]
{Arithmetic Shadows of Chiral Algebras}

\author{Raeez Lorgat}
\address{Perimeter Institute for Theoretical Physics,
 Waterloo, ON N2L 2Y5, Canada}
\email{rlorgat@perimeterinstitute.ca}

\date{}

\begin{abstract}
The scalar genus-$1$ shadow channel of a chiral algebra has Fourier
coefficients $-24\kappa(\cA)\sigma_1(n)$, hence Dirichlet series
\[
D_2(\cA,s)=-24\kappa(\cA)\zeta(s)\zeta(s-1).
\]
For lattice vertex algebras, the constrained Epstein zeta function
on the punctured lattice sector recovers the Riemann zeta function
in rank one and, in higher rank, has completed Mellin transform
given by the lattice theta series. Ordered chiral homology at integer level
recovers the Verlinde formula; the Verlinde dimensions are
polynomials whose leading coefficients are normalized even zeta
values. At genus~$2$, the bar-complex theta component connects to
Siegel modular forms and to the Furusawa--Morimoto form of
B\"ocherer's theorem.
\end{abstract}

\subjclass[2020]{17B69, 11M06, 14H10, 11F46, 81T40}

\keywords{Chiral algebras, shadow obstruction tower, modular
characteristic, Eisenstein series, Verlinde formula, Siegel modular
forms, Epstein zeta function}

\maketitle

\setcounter{tocdepth}{2}
\tableofcontents

% ================================================================
% SECTION 1
% ================================================================
\section{Introduction}\label{sec:intro}

\subsection{The arithmetic content of the shadow tower}

The bar complex of a chiral algebra $\cA$ on an algebraic curve~$X$
is a cochain complex
$\barB^{\ord}(\cA) = T^c(s^{-1}\bar\cA)$
with deconcatenation coproduct and differential encoding the OPE
through the logarithmic kernel $\eta = d\log(z_1 - z_2)$. Over a
family of genus-$g$ curves, the fiberwise differential squares to
$d_{\mathrm{fib}}^2 = \kappa(\cA) \cdot \omega_g$, where
$\omega_g = c_1(\lambda)$ is the first Chern class of the Hodge
bundle on $\overline{\cM}_g$. The single scalar
$\kappa(\cA)$, the \emph{modular characteristic}, controls the
obstruction to flatness of the bar complex as a family.

The shadow obstruction tower $\{S_r(\cA)\}_{r \ge 2}$ extracts the
degree-$r$ constant terms of the universal Maurer--Cartan element
$\Theta_\cA \in \MC(\gmod)$, with $S_2 = \kappa$. The shadow
coefficients are rational functions of the central charge and level,
carrying algebraic data from the OPE into the tower.

These rational numbers carry arithmetic content beyond their
algebraic origin.

\subsection{The Heisenberg base point}

The Heisenberg algebra $\cH_k$ is the abelian Kac--Moody algebra at
level~$k$, with OPE $J(z)J(w) \sim k/(z-w)^2$ and modular
characteristic $\kappa(\cH_k) = k$.
The genus-$1$ amplitude of $\cH_k$ over the family
$\cC_\tau \to \mathfrak{H}$ is
\begin{equation}\label{eq:heis-genus1}
  F_1^{\conn}(\cH_k; q)
  = -k \sum_{n \ge 1} \sigma_1(n)\, q^n
  = k \cdot \frac{E_2^*(\tau) - 1}{24},
\end{equation}
where $E_2^*(\tau) = 1 - 24\sum_{n \ge 1} \sigma_1(n)\,q^n$ is the
quasi-modular Eisenstein series and $\sigma_1(n) = \sum_{d \mid n} d$
is the sum-of-divisors function. The Fourier coefficients are
divisor sums: the simplest arithmetic functions in analytic number
theory.

The Heisenberg computation is the base model. The affine
Kac--Moody, Virasoro, $\cW_N$, and lattice families add the
Lie-theoretic and arithmetic terms one at a time.

\subsection{Summary of results}

The main arithmetic surfaces are the Eisenstein factorisation of
$D_2$, the constrained Epstein zeta of lattice vertex algebras, the
elementary convergence range for the formal shadow Dirichlet series,
the Verlinde polynomial family, and the genus-$2$ Siegel modular
projection problem.

Section~\ref{sec:bocherer} connects the genus-$2$ Maurer--Cartan
element to Siegel modular forms via the Saito--Kurokawa lift and the
B\"ocherer theorem.

Section~\ref{sec:epstein} develops the constrained Epstein zeta
function, its Hecke decomposition, and the Dirichlet--sewing lift.

\subsection{Conventions}

All chiral algebras are $\bZ$-graded with cohomological convention
($|d| = +1$). The modular characteristic is
\[
\kappa(\cA) :=
\begin{cases}
k & \text{Heisenberg } \cH_k, \\[3pt]
\displaystyle\frac{\dim(\fg)(k + h^\vee)}{2h^\vee}
  & \text{affine Kac--Moody } V_k(\fg), \\[6pt]
c/2 & \text{Virasoro } \Vir_c, \\[3pt]
c\,(H_N - 1) & \text{principal } \cW_N,
\end{cases}
\]
where $H_N = \sum_{j=1}^N 1/j$ is the $N$-th harmonic number.
We write $q = e^{2\pi i \tau}$ for the nome. The ordered bar complex
is $\barB^{\ord}(\cA) = T^c(s^{-1}\bar\cA)$, where
$\bar\cA = \ker(\varepsilon)$ is the augmentation ideal and
$|s^{-1}v| = |v| - 1$ is the desuspension.

\subsection{Relation to prior work}

The modular-form techniques used in
Sections~\ref{sec:shadow-eisenstein}--\ref{sec:bocherer} draw
on the classical theory of Eisenstein series, Hecke operators,
and $L$-functions as developed in Zagier~\cite{Zagier08} and
the systematic treatment of analytic methods for $L$-functions
and automorphic forms in Iwaniec--Kowalski~\cite{IK04}.
The factorisation of the shadow Dirichlet series
$D_2(\cA,s)$ into products of Riemann zeta functions
(Theorem~\ref{thm:shadow-eisenstein-standalone}) is an instance of
the Rankin--Selberg method, connecting the arithmetic content
of OPE data to classical multiplicative number theory.
At higher rank, the constrained Epstein zeta function is the
Mellin transform of a theta series and decomposes through the
Hecke eigenbasis of the corresponding modular-form space.

% ================================================================
% SECTION 2
% ================================================================
\section{The shadow Eisenstein theorem}\label{sec:shadow-eisenstein}

\subsection{Genus-$1$ amplitudes from the bar complex}

Let $\cA$ be a chirally Koszul algebra with modular characteristic
$\kappa = \kappa(\cA)$. In the scalar shadow channel, or in a
diagonal Heisenberg/weight-multiset model satisfying the required
Hilbert-Schmidt sewing estimates, the genus-$1$ connected free
energy has Fourier coefficients proportional to $\sigma_1(n)$. The
full genus-$1$ determinant requires the separate analytic sewing
hypotheses.

The genus-$1$ shadow amplitude is
\begin{equation}\label{eq:Sh2}
  \mathrm{Sh}_2^{(1)}(\cA, \tau)
  = \kappa(\cA) \cdot E_2^*(\tau),
\end{equation}
where $E_2^*(\tau) = 1 - 24\sum_{n \ge 1}\sigma_1(n)\,q^n$ is the
quasi-modular Eisenstein series (weight~$2$, with the non-holomorphic
completion $\hat E_2(\tau) = E_2^*(\tau) - 3/(\pi\,\mathrm{Im}(\tau))$
absorbing the anomaly from $d\log E(z,w)$). The Fourier coefficients
$\sigma_1(n) = \sum_{d \mid n} d$ are the sum-of-divisors function.

\begin{example}[Explicit genus-$1$ amplitudes]\label{ex:genus1-families}
  For the standard families:
  \begin{enumerate}[label=\textup{(\alph*)}]
  \item \emph{Heisenberg $\cH_k$:}
    $\kappa(\cH_k) = k$;
    $\mathrm{Sh}_2^{(1)} = k \cdot E_2^*(\tau)$.
  \item \emph{Affine Kac--Moody $V_k(\fg)$:}
    $\kappa(V_k(\fg)) = \dim(\fg)(k + h^\vee)/(2h^\vee)$;
    $\mathrm{Sh}_2^{(1)} = \dim(\fg)(k{+}h^\vee)/(2h^\vee)
    \cdot E_2^*(\tau)$.
    At $k = 0$, $\kappa = \dim(\fg)/2$, so the Heisenberg summand
    contributes a nonzero amplitude.
    At $k = -h^\vee$: $\kappa = 0$; the amplitude vanishes (critical
    level).
  \item \emph{Virasoro $\Vir_c$:}
    $\kappa(\Vir_c) = c/2$;
    $\mathrm{Sh}_2^{(1)} = (c/2) \cdot E_2^*(\tau)$.
    Self-dual at $c = 13$, where $\kappa + \kappa' = 13$.
  \end{enumerate}
\end{example}

\subsection{The Dirichlet series of divisor sums}

The Fourier coefficients $\sigma_1(n)$ assemble into a Dirichlet
series with a classical factorisation.

\begin{theorem}[Shadow Eisenstein theorem]\label{thm:shadow-eisenstein-standalone}
  Let $\cA$ be a chirally Koszul algebra with modular
	  characteristic $\kappa(\cA)$. Define
	  $a_n(\cA)=-24\kappa(\cA)\sigma_1(n)$. The scaled Dirichlet series
	  \begin{equation}\label{eq:D2}
	    D_2(\cA, s)
	    := \sum_{n \ge 1} a_n(\cA)\, n^{-s}
	  \end{equation}
	  satisfies
	  \begin{equation}\label{eq:D2-factored}
	    D_2(\cA, s)
	    = -24\kappa(\cA)\,\zeta(s)\,\zeta(s - 1),
  \end{equation}
  where $\zeta$ is the Riemann zeta function. The genus-$1$
  amplitude Dirichlet series is therefore
  \begin{equation}\label{eq:genus1-dirichlet}
    -24\,\kappa(\cA) \cdot \zeta(s) \cdot \zeta(s - 1):
  \end{equation}
  an Eisenstein $L$-function with the modular characteristic as the
  sole chiral-algebra input.
\end{theorem}

\begin{proof}
  The identity $\sigma_1(n) = \sum_{d \mid n} d$ is a Dirichlet
  convolution: $\sigma_1 = \id * \mathbf{1}$, where
  $\id(n) = n$ and $\mathbf{1}(n) = 1$. The Dirichlet series of a
  convolution is the product of the constituent Dirichlet series:
  \[
    \sum_{n \ge 1} \sigma_1(n)\, n^{-s}
    = \Bigl(\sum_{n \ge 1} n^{1-s}\Bigr)
      \Bigl(\sum_{n \ge 1} n^{-s}\Bigr)
    = \zeta(s - 1)\,\zeta(s).
  \]
  The factor $-24\,\kappa(\cA)$ comes from the
  relation~\eqref{eq:Sh2}: the genus-$1$ amplitude has Fourier
  coefficient $-24\,\kappa(\cA)\,\sigma_1(n)$ at $q^n$.
\end{proof}

\begin{remark}[Analytic properties]\label{rem:D2-analytic}
  The product $\zeta(s)\,\zeta(s-1)$ has simple poles at $s = 1$
  and $s = 2$, a functional equation inherited from the individual
  zeta factors, and zeros at the nontrivial zeros of $\zeta$ shifted
  by $0$ and $1$. The residue at $s = 2$ is
  $\mathrm{Res}_{s=2}\,\zeta(s)\zeta(s-1)=\zeta(2)=\pi^2/6$.
\end{remark}

\begin{remark}[Non-automorphicity of the shadow tower]
  \label{rem:non-automorphic}
  The shadow $L$-function
  $L^{\sh}_\cA(s) = \sum_{r \ge 2} S_r\, r^{-s}$, built from the
  shadow tower coefficients $S_r$, is a different object from
  $D_2(\cA,s)$. No Euler product or Selberg-class structure is part
  of its definition. Analytic continuation requires an explicit
  generating-function theorem for the full shadow tower.
\end{remark}

\subsection{The Polyakov correspondence}

The shadow Eisenstein theorem has a physical counterpart.
Polyakov's determinant calculation computes the string partition function by
gauge-fixing the worldsheet metric to conformal gauge, producing a
conformal anomaly $\langle T^a{}_a \rangle = (c/12)\,R$ and a ghost
system with $c_{\mathrm{ghost}} = -26$. The bar-cobar side records
the corresponding shadow coefficients.

\begin{center}
\renewcommand{\arraystretch}{1.15}
\begin{tabular}{@{} l l @{}}
\toprule
\textbf{Polyakov} & \textbf{Bar-cobar} \\ \midrule
Conformal anomaly $(c/12)\,R$
  & Modular characteristic $\kappa(\cA)$ \\
$\det'_\zeta \Delta$ (Laplacian)
  & Fredholm $\det(1 - K)$ (Bergman kernel) \\
Critical dimension $c = 26$
  & $\kappa(\Vir_{26-c}) = 0$ \\
Ghost system $bc$
  & Koszul dual $\cA^!$ \\
$F_g = \kappa\,\lambda_g^{\FP}$ (uniform)
  & Genus-$g$ shadow amplitude \\
\bottomrule
\end{tabular}
\end{center}

The genus-$1$ functional determinant
$\det'_\zeta \Delta = |\eta(\tau)|^4\,(\mathrm{Im}\,\tau)^2$ is the
Heisenberg contribution; the divisor sums in $E_2^*$ encode the
one-loop corrections. The shadow Eisenstein theorem extracts the
multiplicative structure of these corrections.


% ================================================================
% SECTION 3
% ================================================================
\section{The categorical zeta function}\label{sec:categorical-zeta}

\subsection{Lattice vertex algebras and the punctured lattice sector}

Let $\Lambda$ be an even lattice of rank~$r$ with inner product
$\langle\cdot,\cdot\rangle$, and let $V_\Lambda$ be the associated
lattice vertex algebra. The vertex algebra $V_\Lambda$ decomposes as
$V_\Lambda = \cH^{\otimes r} \otimes \bC[\Lambda]$, where
$\cH^{\otimes r}$ is the rank-$r$ Heisenberg algebra and
$\bC[\Lambda]$ is the group algebra. The modular characteristic is
$\kappa(V_\Lambda) = r$.

The augmentation removes the vacuum vector. The lattice-sector
Epstein sum therefore runs over the punctured lattice
$\Lambda\setminus\{0\}$; this set is not a sublattice.

\begin{definition}[Constrained Epstein zeta function]
  \label{def:constrained-epstein}
  For a lattice vertex algebra $V_\Lambda$ of rank~$r$, the
  \emph{constrained Epstein zeta function} is
  \begin{equation}\label{eq:constrained-epstein}
    \varepsilon^c_s(V_\Lambda)
    :=
    {\sideset{}{'}\sum}_{\ell \in \Lambda}
    \|\ell\|^{-2s}
    = \sum_{\substack{\ell \in \Lambda \\ \ell \ne 0}}
    \langle \ell, \ell \rangle^{-s},
  \end{equation}
  where the prime denotes omission of $\ell = 0$.
\end{definition}

This is the Epstein zeta function of the quadratic form
$Q(\ell) = \langle \ell, \ell \rangle$ with the zero vector omitted.
For a general even lattice, the Epstein zeta function
$E_\Lambda(s) = \sum_{\ell \ne 0} Q(\ell)^{-s}$ is a classical
object in analytic number theory, studied by Epstein~\cite{Epstein03,
Epstein07} and Siegel~\cite{Siegel51}. The constraint to
$\Lambda^{\mathrm{bar}}$ is the bar-complex restriction: the
augmentation ideal removes the vacuum, and the surviving lattice
vectors carry the OPE data.

\subsection{Recovery of the Riemann zeta function}

\begin{theorem}[Rank-$1$ recovery]\label{thm:rank1-recovery}
  For the even rank-$1$ lattice $A_1$ with quadratic form
  $Q(m)=2m^2$, the constrained Epstein zeta function is
  \begin{equation}\label{eq:rank1-epstein}
    \varepsilon^c_s(V_{A_1})
    = 2 \cdot (2)^{-s} \sum_{m \ge 1} m^{-2s}
    = 2^{1-s}\,\zeta(2s).
  \end{equation}
  With the unit quadratic form $Q(m)=m^2$ one obtains
  $2\zeta(2s)$. The rank-one Narain momentum-winding lattice at
  radius $R$ gives
  $2(R^{2s}+R^{-2s})\zeta(2s)$, hence $4\zeta(2s)$ at $R=1$.
\end{theorem}

\begin{proof}
  At rank~$1$ with $\Lambda = \bZ$ and quadratic form
  $Q(m) = m^2$ (unit radius), the nonzero lattice vectors are
  $\pm 1, \pm 2, \ldots$, each contributing $|m|^{-2s}$. By
  symmetry $m \mapsto -m$:
  \[
    \varepsilon^c_s(V_\bZ)
    = 2 \sum_{m=1}^\infty m^{-2s}
    = 2\,\zeta(2s).
  \]
  For the even lattice with $Q(m) = 2m^2$:
  $\varepsilon^c_s = 2 \cdot (2)^{-s}\,\zeta(2s) =
  2^{1-s}\zeta(2s)$. The Narain formula counts the two
  momentum-winding directions.
\end{proof}

\begin{remark}[Functional equation]\label{rem:epstein-fe}
  The Epstein zeta function of a rank-$r$ lattice $\Lambda$ satisfies
  the functional equation
  \[
    \pi^{-s}\,\Gamma(s)\,E_\Lambda(s)
    = |\det G|^{-1/2}\,
    \pi^{-(r/2-s)}\,\Gamma(r/2-s)\,E_{\Lambda^*}(r/2-s),
  \]
  where $G$ is the Gram matrix of $\Lambda$ and $\Lambda^*$ is the
  dual lattice. For the self-dual lattice $\Lambda = \Lambda^*$ at
  rank~$1$, this reduces to the functional equation of
  $\zeta(2s)$. The constrained Epstein zeta inherits this functional
  equation from the standard Epstein zeta.
\end{remark}

\subsection{Higher-rank decomposition}

At rank $r \ge 2$, the Epstein zeta function is no longer a product
of Riemann zeta functions. The Hecke decomposition provides the
correct generalisation.

\begin{proposition}[Theta-series Hecke decomposition]
  \label{prop:hecke-decomposition}
  For a lattice vertex algebra $V_\Lambda$ of rank~$r$ with
  Gram matrix~$G$, the completed constrained Epstein zeta is the
  Mellin transform of $\Theta_\Lambda-1$. Decomposing the theta
  series into a Hecke eigenbasis gives
  \begin{equation}\label{eq:hecke-decomp}
    \pi^{-s}\Gamma(s)\varepsilon^c_s(V_\Lambda)
    = \sum_j c_j\,\Lambda(s,f_j),
  \end{equation}
  where the $f_j$ range over Eisenstein and cuspidal Hecke
  eigenforms occurring in the modular-form space generated by
  $\Theta_\Lambda$, and $\Lambda(s,f_j)$ are their completed
  $L$-functions.
\end{proposition}

\begin{proof}
  The theta series $\Theta_\Lambda$ is a modular form of weight
  $r/2$ for the appropriate congruence subgroup. The finite-dimensional
  modular-form space containing it decomposes into Eisenstein and
  cuspidal Hecke eigenspaces. The Mellin transform of
  $\Theta_\Lambda(iy)-1$ is the Epstein zeta integral, and Mellin
  transform sends each Hecke eigenform summand to its completed
  $L$-function.
\end{proof}

\begin{remark}[Lattice theta function]
  \label{rem:lattice-theta}
  The lattice theta function
  $\Theta_\Lambda(\tau) = \sum_{\ell \in \Lambda}
  q^{Q(\ell)/2}$ is a modular form for a congruence subgroup
  of~$\mathrm{SL}_2(\bZ)$, and the constrained Epstein zeta is
  its Mellin transform:
  \[
  \pi^{-s}\Gamma(s)\,\varepsilon^c_s(V_\Lambda)
  =
  \int_0^\infty(\Theta_\Lambda(iy)-1)y^{s-1}\,dy .
  \]
  The subtraction of~$1$ removes the zero vector.
\end{remark}

The Niemeier lattice vertex algebras $V_\Lambda$ at rank~$24$ have
modular characteristic $\kappa = 24$, independent of the root
system. The lattice theta function discriminates between the
$24$ Niemeier lattices; the shadow tower, which depends only
on~$\kappa$, does not. The constrained Epstein zeta function
$\varepsilon^c_s(V_\Lambda)$ carries this finer arithmetic.


% ================================================================
% SECTION 4
% ================================================================
\section{Depth decomposition of the shadow $L$-function}
\label{sec:depth-decomposition}

\subsection{The shadow $L$-function}

The shadow obstruction tower $\{S_r(\cA)\}_{r \ge 2}$ is the
sequence of degree-$r$ projections of the Maurer--Cartan element
$\Theta_\cA \in \MC(\gmod)$, with $S_2 = \kappa(\cA)$.

\begin{definition}[Shadow $L$-function]\label{def:shadow-L}
  For a chirally Koszul algebra $\cA$, the \emph{shadow $L$-function}
  is the Dirichlet series
  \begin{equation}\label{eq:shadow-L}
    L^{\sh}_\cA(s)
    := \sum_{r \ge 2} S_r(\cA)\, r^{-s}.
  \end{equation}
\end{definition}

\begin{example}[Shadow $L$-functions by family]
  \label{ex:shadow-L-families}
  The four depth classes produce qualitatively different
  $L$-functions.
  \begin{enumerate}[label=\textup{(\alph*)}]
  \item \emph{Class G} (Heisenberg, $r_{\max} = 2$):
    $S_r = 0$ for $r \ge 3$, so
    $L^{\sh}_{\cH_k}(s) = k \cdot 2^{-s}$.
    Class G gives a one-term Dirichlet polynomial.
  \item \emph{Class L} (affine Kac--Moody, $r_{\max} = 3$):
    $L^{\sh}_{V_k(\fg)}(s) = \kappa \cdot 2^{-s}
    + S_3 \cdot 3^{-s}$.
    Class L gives a two-term Dirichlet polynomial.
  \item \emph{Class C} ($\beta\gamma$, $r_{\max} = 4$):
    Three terms; the quartic shadow coefficient $S_4$ controls the
    discriminant $\Delta = 8\kappa S_4$.
  \item \emph{Class M} (Virasoro, $\cW_N$; $r_{\max} = \infty$):
    The full Dirichlet series is infinite.
    For Virasoro:
    $S_4(\Vir_c) = 10/[c(5c+22)]$.
  \end{enumerate}
\end{example}

\subsection{Analytic properties}

\begin{proposition}[Elementary convergence of
  $L^{\sh}_\cA$]\label{prop:shadow-L-analytic}
  For a class-M algebra with shadow coefficients
  $S_r = O(r^{-\alpha})$ for some $\alpha > 0$, the shadow
  Dirichlet series $L^{\sh}_\cA(s)$ converges absolutely for
  $\Re(s)>1-\alpha$.
\end{proposition}

\begin{proof}
  Absolute convergence follows from comparison with
  $\sum r^{-\alpha-\Re(s)}$.
\end{proof}

No Euler product or Selberg-class structure is part of the definition
of $L^{\sh}_\cA$. Meromorphic continuation, poles, and special values
require a separate admissible generating-function theorem for the
full shadow tower.

\subsection{Depth stratification}

The four depth classes $\mathbf{G}/\mathbf{L}/\mathbf{C}/\mathbf{M}$
induce a filtration on the space of shadow $L$-functions.

\begin{definition}[Depth filtration]\label{def:depth-filtration}
  The \emph{depth-$d$ truncation} of the shadow $L$-function is
  \begin{equation}\label{eq:depth-truncation}
    L^{\sh,\le d}_\cA(s)
    := \sum_{r=2}^{d} S_r(\cA)\, r^{-s}.
  \end{equation}
  The \emph{depth defect} is
  $\delta L^{\sh}_\cA(s)
  := L^{\sh}_\cA(s) - L^{\sh,\le d}_\cA(s)
  = \sum_{r > d} S_r\, r^{-s}$.
\end{definition}

For class G: $\delta L^{\sh, \le 2}_\cA = 0$ (no defect).
For class L: $\delta L^{\sh, \le 3}_\cA = 0$.
For class C: $\delta L^{\sh, \le 4}_\cA = 0$.
For class M: $\delta L^{\sh, \le d}_\cA \ne 0$ for all
finite~$d$; the defect is genuinely infinite.

The scalar discriminant $\Delta(\cA)=8\kappa(\cA)S_4(\cA)$ detects
the quartic/contact obstruction. It does not distinguish class C
from class M. The coarse criteria are:
\[
\mathbf G:\ S_3=S_4=0,\qquad
\mathbf L:\ S_3\ne0,\ S_4=0,\qquad
\mathbf C:\ S_3=0,\ S_4\ne0,\ S_{r\ge5}=0,
\]
while class $\mathbf M$ is nonterminating.


% ================================================================
% SECTION 4B: ARITHMETIC DUALITY
% ================================================================
\section{Arithmetic duality of Bernoulli primes and characteristic
 primes}\label{sec:arithmetic-duality}

This section anchors the arithmetic duality between the Verlinde
polynomial family $\{Z_g(k)\}$ and the shadow tower
$\{S_r(\mathrm{Vir}_c)\}$. The comparison uses Bernoulli
numerators on the Verlinde side and cleared weighted Riccati
numerators on the Virasoro shadow side.

\subsection{The duality}

\begin{proposition}[Finite $Z_g$--$S_r$ arithmetic separation]
\label{prop:z-g-s-r-arithmetic-duality-standalone}
The leading Kummer-irregular primes $\{691, 3617\}$ are present on
$Z_g$ through the Hurwitz--Bernoulli leading term $B_{2g-2}$ and are
absent from the cleared weighted Virasoro shadow numerators
$\widehat N_r$ through $r = 11$. Explicitly:
\begin{enumerate}[label=\textup{(\roman*)}]
\item\label{duality-i}
 $B_{12} = -691/2730$ contributes $691 \mid \mathrm{num}(Z_7)$.
\item\label{duality-ii}
 $B_{16} = -3617/510$ contributes
 $3617 \mid \mathrm{num}(Z_9)$.
\item\label{duality-iii}
 $\{691, 3617\}\cap \mathrm{supp}(\widehat N_r)=\varnothing$
 for $2 \le r \le 11$ in the computed weighted Riccati range.
\end{enumerate}
\end{proposition}

\begin{proof}
The first two assertions are the displayed Bernoulli values. The
third assertion is the finite factorisation of the cleared weighted
Riccati numerators $\widehat N_r$ for $2\le r\le 11$.
\end{proof}

The finite comparison is sharp only for the displayed verification
window. The prime $2111$ appears in the Riccati numerator
factorisation $29554 = 2 \cdot 7 \cdot 2111$, but it is not
witnessed on the Verlinde leading term for $g\le9$ because
$2111\nmid\mathrm{num}(B_{2m})$ for $2m\le16$. This is not an
ordering statement about Kummer-irregular primes.

\subsection{Riccati-arithmetic characteristic primes}

The cleared weighted numerators $\widehat N_r(c)$ through $r=11$
contain the Riccati-propagated prime factors
$\{61,193,2111,16657,3988097\}$. These primes occur in the integer
coefficients of the primitive-cleared polynomials obtained from the
Virasoro shadow recursion $H(t)^2=t^4Q_L(t)$.

\begin{remark}[Kummer and Riccati-prime separation]
\label{rem:kummer-riccati-prime-separation}
In the computed range $2\le r\le 11$, the primes
$\{1423, 23, 43, 419\}$ are absent from the Kummer-irregular list
and occur only on the Riccati side. The prime $3067$ is not witnessed
as Kummer-irregular at the current verification depth. The structural
distinction is geometric: Kummer-arithmetic primes track $B_{2g-2}$
on the Verlinde side, while Riccati-arithmetic primes track the Kac
determinant on the shadow side.
\end{remark}

\begin{remark}[Complement: $S_r$-characteristic primes, absent from
leading $Z_g$]\label{rem:sr-characteristic-primes}
The characteristic primes $\{61, 193, 2111, 16657, 3988097\}$
occur in $\mathrm{num}(S_r(\mathrm{Vir}_c))$ for suitable $r$ but
do \emph{not} divide $\mathrm{num}(B_{2g-2})$ at $g \le 9$. Their
emergence is Riccati-seed-driven: they occur in the integer
coefficients produced by the shadow recursion and are invisible to
the Hurwitz--Bernoulli leading term on $Z_g$ in the displayed range.
\end{remark}

% ================================================================
% SECTION 5
% ================================================================
\section{Verlinde polynomials from ordered chiral homology}
\label{sec:verlinde}

\subsection{Recovery of the Verlinde formula}

At integer level $k \ge 1$, the quantum group parameter
$q = e^{2\pi i/(k+2)}$ is a root of unity, the representation
category of $\widehat{\fsl}_2$ at level~$k$ truncates to $k+1$
integrable modules, and the ordered chiral homology at each genus
computes a finite-dimensional invariant. Write $n = k + 2$ for the
shifted level.

\begin{proposition}[Verlinde formula from ordered chiral
  homology]\label{prop:verlinde-recovery}
  Let $k \ge 1$ be a positive integer, and let
  $S_{jl} = \sqrt{2/n}\,\sin\!\bigl(\pi(j+1)(l+1)/n\bigr)$
  be the modular $S$-matrix for $\widehat{\fsl}_2$ at level~$k$,
  where $j, l \in \{0, 1, \ldots, k\}$.
  \begin{enumerate}[label=\textup{(\roman*)}]
  \item \textup{(Genus~$0$.)}
    $\dim H^0\!\bigl(\cM_{0,0},\, V_k(\fsl_2)\bigr)
    = \sum_{j=0}^{k} S_{0j}^2 = 1$.
  \item \textup{(Genus~$1$.)}
    $\dim H^0\!\bigl(\cM_{1,0},\, V_k(\fsl_2)\bigr)
    = \sum_{j=0}^{k} S_{0j}^0 = k + 1$.
  \item \textup{(General genus.)}
    $Z_g(k) := \dim H^0\!\bigl(\cM_{g,0},\,
    V_k(\fsl_2)\bigr)
    = \sum_{j=0}^{k} S_{0j}^{2-2g}$.
  \item \textup{(Handle attachment.)}
    $Z_{g+1} = \sum_{j=0}^{k} S_{0j}^{-2} \cdot S_{0j}^{2-2g}
    = \sum_{j=0}^{k} S_{0j}^{2-2(g+1)}$.
  \item \textup{(Separating factorization.)}
    Under $\Sigma_g \rightsquigarrow \Sigma_{g_1} \cup \Sigma_{g_2}$
    with $g = g_1 + g_2$:
    $Z_g = \sum_{j=0}^{k} S_{0j}^{2-2g_1} \cdot S_{0j}^{2-2g_2}
    \cdot S_{0j}^{-2}$.
  \end{enumerate}
\end{proposition}

\begin{proof}
  The identity $\sum_j S_{0j}^2 = 1$ is the $(0,0)$-entry of
  $S \cdot S^T = I$, following from the orthogonality relation
  $\sum_{j=0}^{k} \sin(\pi(j{+}1)a/n)
  \sin(\pi(j{+}1)b/n) = n\,\delta_{ab}/2$.
  Since $S_{0j}^0 = 1$ for all~$j$, $Z_1 = k + 1$.
  The KZB local system at integrable level decomposes
  into $k+1$ channels labeled by integrable representations; each
  channel~$j$ contributes $S_{0j}^{2-2g}$ to the partition function.
  Handle attachment and separating factorization are the same
  diagonal channel computation under sewing.
\end{proof}

The first values are:
\begin{center}
\renewcommand{\arraystretch}{1.3}
\begin{tabular}{@{} c r r r r r @{}}
\toprule
$k$ & $Z_0$ & $Z_1$ & $Z_2$ & $Z_3$ & $Z_4$ \\
\midrule
$1$ & $1$ & $2$ & $4$ & $8$ & $16$ \\
$2$ & $1$ & $3$ & $10$ & $36$ & $136$ \\
$3$ & $1$ & $4$ & $20$ & $120$ & $800$ \\
$4$ & $1$ & $5$ & $35$ & $329$ & $3{,}611$ \\
$5$ & $1$ & $6$ & $56$ & $784$ & $13{,}328$ \\
\bottomrule
\end{tabular}
\end{center}

\begin{remark}[Two genus invariants]\label{rem:Fg-vs-Zg}
  The shadow coefficients $F_g$ and the Verlinde dimensions $Z_g$
  are different invariants of the same algebra. The shadow
  coefficient $F_g = \kappa \cdot \lambda_g^{\FP}$
  varies continuously with the level and is a small rational
  number. For $V_1(\fsl_2)$ one has $\kappa=9/4$, hence
  $F_2 = 7/2560$. The Verlinde
  dimension $Z_g$ is a positive integer defined only at integer
  level, with $Z_2 = 4$ at $k = 1$.
\end{remark}

\subsection{The Verlinde polynomial family}

The Verlinde dimensions at each genus are polynomials in $n = k+2$.

\begin{theorem}[Verlinde polynomial family]
  \label{thm:verlinde-polynomial}
  For $\widehat{\fsl}_2$ at level $k \ge 1$ and genus $g\ge2$, the Verlinde
  dimension $Z_g(k) = \sum_{j=0}^{k} S_{0j}^{2-2g}$ is a
  polynomial $P_g(n)$ of degree $3(g-1)$ in $n = k+2$ with the
  following structure.
  \begin{enumerate}[label=\textup{(\roman*)}]
  \item \textup{(Universal factor.)} For $g \ge 2$,
    \begin{equation}\label{eq:verlinde-factor}
      P_g(n) = n^{g-1}(n^2 - 1) \cdot R_{g-2}(n^2),
    \end{equation}
    where $R_{g-2}$ is a polynomial of degree $g - 2$ in $n^2$
    with rational coefficients and positive leading term.
  \item \textup{(Explicit formulas through $g = 6$.)}
    \begin{align}
      P_2(n) &= \frac{n(n^2 - 1)}{6},
        \label{eq:P2} \\
      P_3(n) &= \frac{n^2(n^2 - 1)(n^2 + 11)}{180},
        \label{eq:P3} \\
      P_4(n) &= \frac{n^3(n^2 - 1)(2n^4 + 23n^2 + 191)}{7560},
        \label{eq:P4} \\
      P_5(n) &= \frac{n^4(n^2 - 1)(n^2 + 11)
        (3n^4 + 10n^2 + 227)}{226800},
        \label{eq:P5} \\
      P_6(n) &= \frac{n^5(n^2 - 1)(2n^8 + 35n^6
        + 321n^4 + 2125n^2 + 14797)}{2993760}.
        \label{eq:P6}
    \end{align}
  \item \textup{(Leading asymptotics.)} The leading coefficient
    as $n \to \infty$ is
    \begin{equation}\label{eq:verlinde-leading}
      P_g(n) \sim
      \frac{\zeta(2g - 2)}{2^{g-2}\,\pi^{2g-2}}\;
      n^{3(g-1)}
      \qquad (n \to \infty).
    \end{equation}
  \item \textup{(Rational generating function.)} At fixed~$n$,
    \begin{equation}\label{eq:verlinde-gf}
      G_n(x) := \sum_{g \ge 1} Z_g(k)\, x^{g-1}
      = \sum_{j=1}^{n-1}
      \frac{1}{1 - a_j\, x},
      \qquad
      a_j = \frac{n}{2\sin^2(\pi j/n)}.
    \end{equation}
  \end{enumerate}
\end{theorem}

\begin{proof}
  The cosecant-power representation
  $Z_g(k) = (n/2)^{g-1}\sum_{j=1}^{n-1}
  \csc^{2g-2}(\pi j/n)$
  follows from $S_{0j} = \sqrt{2/n}\,\sin(\pi j/n)$. The
  cosecant power sum
  $\sigma_m(n) = \sum_{j=1}^{n-1}\csc^{2m}(\pi j/n)$
  is a polynomial of degree~$2m$ in~$n$ with rational coefficients:
  a classical result of Euler, with leading coefficient
  $\sigma_m(n) \sim 2\zeta(2m)\pi^{-2m}n^{2m}$.
  Parts~(i)--(iii) follow. Part~(iv): the geometric series
  $\sum_{g \ge 1} a_j^{g-1} x^{g-1} = 1/(1 - a_j x)$ applied
  term by term gives the rational generating function. The explicit
  formulas~\eqref{eq:P2}--\eqref{eq:P6} are obtained by polynomial
  interpolation and verified by direct $S$-matrix summation.
\end{proof}

\begin{remark}[Zeta values in the Verlinde tower]
  \label{rem:zeta-values}
  The leading coefficient
  $[n^{3g-3}]P_g = \zeta(2g-2)/(2^{g-2}\pi^{2g-2})$ connects the
  Verlinde tower to zeta values at even positive integers:
  \[
    \frac{\zeta(2)}{2^0\,\pi^2} = \frac{1}{6},
    \quad
    \frac{\zeta(4)}{2\,\pi^4} = \frac{1}{180},
    \quad
    \frac{\zeta(6)}{4\,\pi^6} = \frac{1}{3780},
    \quad
    \frac{\zeta(8)}{8\,\pi^8} = \frac{1}{75600}.
  \]
  The displayed denominators are the reduced leading-coefficient
  denominators of $P_2,P_3,P_4,P_5$.
\end{remark}

\begin{remark}[Structural observations]
  \label{rem:verlinde-structure}
  Three features of the polynomial
  family~\eqref{eq:verlinde-factor}:
  \begin{enumerate}[label=\textup{(\alph*)}]
  \item The factor $n^2 - 1 = (k+1)(k+3)$ vanishes at
    $k = -1$ and $k = -3$, reflecting the boundary of the
    integrable range. The factor $n^{g-1}$ vanishes at $n = 0$
    ($k = -2$, the critical level), where the Verlinde formula
    degenerates.
  \item $P_2(n) = \binom{n+1}{3}$ is the unique genus at which
    $Z_g$ is a single binomial coefficient (the tetrahedral
    numbers).
  \item The factor $(n^2 + 11)$ appears in both $P_3$ and $P_5$,
    suggesting a recurrence in the polynomials $R_{g-2}$.
  \end{enumerate}
\end{remark}


% ================================================================
% SECTION 6
% ================================================================
\section{The B\"ocherer bridge: genus-$2$ Siegel modular forms}
\label{sec:bocherer}

\subsection{From the Maurer--Cartan element to Siegel modular forms}

For lattice VOAs and rational conformal-block systems satisfying the
analytic sewing hypotheses, the genus-$g$ theta or conformal-block
amplitudes transform under $\Sp(2g,\bZ)$. The genus-$2$ lattice case
is the first place where the bar-complex theta component meets the
arithmetic of Siegel modular forms.

The genus-$2$ bar complex on $\Sigma_2$ with Siegel period matrix
$\Omega \in \HHH_2 = \{Z \in M_2(\bC) : Z = Z^t,\;
\mathrm{Im}\,Z > 0\}$ involves the prime form $E(z,w)$ on
$\Sigma_2$ and the matrix $(\mathrm{Im}\,\Omega)^{-1}$ (a $2 \times 2$
matrix, not a scalar; at $g = 1$ this reduces to
$1/\mathrm{Im}(\tau)$). The bar propagator is
$\eta^{(2)}_{ij} = d\log E(z_i, z_j)$, weight~$1$.

\subsection{The Hecke decomposition at genus~$2$}

The Hecke operators $T_p$ act on the space $M_\cA$ of activated
amplitudes, decomposing
$M_\cA = \bigoplus_\chi M_\chi$
into generalised eigenspaces. The \emph{arithmetic packet connection}
\begin{equation}\label{eq:arith-packet}
  \nabla^{\mathrm{arith}}_\cA
  = d - \bigoplus_\chi
  \frac{\Lambda'_\chi(s)}{\Lambda_\chi(s)}
  (\Id_{M_\chi} + N_\chi)\,ds
\end{equation}
is a flat meromorphic connection on $M_\cA \otimes \cO_\bC$ over
the spectral $s$-line, where $\Lambda_\chi(s)$ is the completed
$L$-function of the eigenform and $N_\chi$ is the nilpotent part.
The singular divisor depends only on the semisimple part
$M^{\mathrm{ss}}_\cA$: the arithmetic is semisimple, the homotopy
defect is unipotent.

\subsection{The Saito--Kurokawa lift and the B\"ocherer theorem}

For even unimodular lattice vertex algebras, the genus-$2$ amplitude
is a Siegel modular form for $\Sp(4,\bZ)$, and B\"ocherer's theorem
relates its Fourier coefficients to central $L$-values in the
Saito--Kurokawa case.

\begin{theorem}[Leech $\chi_{12}$ projection]
  \label{thm:leech-chi12-projection-standalone}
  For the Leech lattice vertex algebra, the genus-$2$ Siegel theta
  component has nonzero projection to the one-dimensional cusp space
  $S_{12}(\Sp(4,\bZ))=\bC\chi_{12}$. The cusp form is
  \[
    \chi_{12}=\SK(f_{22}),
  \]
  where $f_{22}=\Delta E_{10}$ is the weight-$22$ Hecke eigenform.
  For fundamental $D<0$, the Saito--Kurokawa/B\"ocherer coefficient
  $B(D;\chi_{12})$ equals the half-integral-weight coefficient
  $c(|D|)$, and Waldspurger gives
  \begin{equation}\label{eq:waldspurger-leech-standalone}
  |c(|D|)|^2
  \;\propto\;
  L(11,f_{22}\times\chi_D)\,|D|^{10}.
  \end{equation}
\end{theorem}

\begin{proof}
The cusp-space dimension and the identity
$\chi_{12}=\SK(f_{22})$ are the Igusa--Saito--Kurokawa input. The
nonzero projection is computed from the Igusa relation and the
genus-$2$ Leech representation numbers; the coefficient is
$c_2\approx -1.918\times10^{-6}$ with zero residual in the exact
rational computation. The last assertion is Waldspurger's theorem
applied to the half-integral-weight form attached to the
Saito--Kurokawa lift; Furusawa--Morimoto~\cite{FurusawaMorimoto}
give the refined B\"ocherer form.
\end{proof}

\begin{remark}[Arithmetic input]\label{rem:weil-analogy}
Deligne's theorem supplies the Ramanujan bound for the Hecke
eigenforms appearing in the genus-$2$ decomposition. The
Maurer--Cartan equation supplies the chiral source of the theta
component; it does not replace the arithmetic proof of the bound.
\end{remark}


% ================================================================
% SECTION 7
% ================================================================
\section{The constrained Epstein zeta and the Dirichlet--sewing lift}
\label{sec:epstein}

\subsection{The Dirichlet--sewing lift}

The genus-$1$ connected free energy
$F_1^{\conn}(\cA; q) = -\log\det(1 - K_q)$ encodes connected sewing
amplitudes $a_\cA(N)$ as Fourier coefficients. The
\emph{Dirichlet--sewing lift} is
\begin{equation}\label{eq:dirichlet-sewing}
  S_\cA(u)
  = \sum_{N \ge 1} a_\cA(N)\, N^{-u}
  = \frac{(2\pi)^u}{\Gamma(u)}
  \int_0^\infty F_\cA^{\conn}(e^{-2\pi y})\, y^{u-1}\, dy.
\end{equation}
This is a Dirichlet series whose analytic structure is controlled by
$\Theta_\cA$.

\begin{example}[Sewing Dirichlet series by family]
  \label{ex:sewing-dirichlet}
  \begin{align}
    S_{\cH}(u) &= \zeta(u)\,\zeta(u+1),
      \label{eq:sewing-heis} \\
    S_{V_k(\fg)}(u) &= \dim(\fg) \cdot \zeta(u)\,\zeta(u+1),
      \label{eq:sewing-km} \\
    S_{\Vir}(u) &= \zeta(u+1)\,(\zeta(u) - 1),
      \label{eq:sewing-vir} \\
    S_{\cW_N}(u) &= \zeta(u+1)\Bigl((N-1)\,\zeta(u)
      - \sum_{j=1}^{N-1} H_j(u)\Bigr),
      \label{eq:sewing-wn}
  \end{align}
  where $H_j(u) = \sum_{m=1}^j m^{-u}$ are partial harmonic sums.
\end{example}

\begin{remark}[The Heisenberg base point]\label{rem:sewing-heis}
  The Heisenberg sewing Dirichlet series
  $S_\cH(u) = \zeta(u)\,\zeta(u+1)$ is a product of two shifted
  Riemann zeta functions, connecting the sewing amplitudes to
  the convolution structure of divisor sums. Every other family's
  sewing series decomposes as the Heisenberg core plus a defect:
  the affine Kac--Moody series is $\dim(\fg)$ copies of Heisenberg,
  the Virasoro series subtracts the constant term~$1$, and $\cW_N$
  has $N-1$ Heisenberg copies minus partial harmonic corrections.
\end{remark}

\subsection{The three-tier Euler--Koszul classification}

The \emph{Euler defect}
$D_\cA(u) := S_\cA(u)/(|W(\cA)|\,\zeta(u)\,\zeta(u+1))$,
where $|W(\cA)|$ is the number of strong generators, classifies
chiral algebras into three tiers.

\begin{center}
\renewcommand{\arraystretch}{1.15}
\begin{tabular}{@{} l l l @{}}
\toprule
\textbf{Tier} & \textbf{Condition} & \textbf{Families} \\ \midrule
Exact Euler--Koszul & $D_\cA(u) \equiv 1$
  & $\cH$, $V_k(\fg)$, $\beta\gamma$ \\
Finitely defective & $1 - D_\cA(u)$ finite harmonic ratio
  & Virasoro, $\cW_N$ \\
Non-Euler--Koszul & $S_\cA$ decomposes into Hecke $L$-functions
  & Lattice VOAs with $h(D) \ge 2$ \\
\bottomrule
\end{tabular}
\end{center}

The Euler--Koszul class $\mathrm{ek}(\cA) = \max_i(w_i - 1)$,
where $w_i$ are the conformal weights of the strong generators,
and the shadow depth $r_{\max}(\cA)$ are independent invariants:
character arithmetic and OPE interaction are logically independent
axes of complexity.

\subsection{The two-variable $L$-object}

Under the diagonal weight-multiset sewing hypotheses, Rankin--Selberg
unfolding against $E^*(\tau, s)$ gives
\begin{equation}\label{eq:two-variable-L}
  L_\cA(s, u)
  = \int_{\cF^{\mathrm{reg}}}
  F_\cA^{\conn}(\tau)\, E^*(\tau, s)\, y^{u-2}\, dx\, dy,
\end{equation}
where $\cF^{\mathrm{reg}}$ is the regularised fundamental domain.
The same hypotheses give the one-variable reduction
\begin{equation}\label{eq:sewing-hecke}
  L_\cA(s, u)
  = \frac{\Gamma(v)}{(2\pi)^v}\, S_\cA(v),
  \qquad
  v = s + u - 1.
\end{equation}
The scattering poles of $E^*$ at $s = \rho/2$ for nontrivial zeros
$\rho$ of $\zeta$ become structural obstructions to an Euler product
for the two-variable object.

\subsection{Constrained Epstein zeta at higher rank}

For a lattice VOA $V_\Lambda$, the algebraic shadow depth and the
arithmetic Epstein depth are separate invariants. In the diagonal
sector the critical-line count has the form
\[
d=1+d_{\mathrm{arith}}+d_{\mathrm{alg}},
\]
where $d_{\mathrm{arith}}$ comes from the automorphic components of
the theta series and $d_{\mathrm{alg}}$ records algebraic
finite-defect factors. Lattice VOAs remain Gaussian in the shadow
depth sense; the Epstein arithmetic depth records the lattice rank
and theta decomposition.

\begin{remark}[The finite Miura defect]\label{rem:miura-defect}
  The principal $\cW_N$ character is
  $\chi_{\cW_N}(q)
  = \prod_{n \ge 2}(1 - q^n)^{-\min(n-1, N-1)}$;
  the Heisenberg character is
  $\chi_\cH(q) = \prod_{n \ge 1}(1 - q^n)^{-1}$. The ratio
  \[
    \chi_{\cW_N}(q)
    = \chi_\cH(q)^{N-1}
    \cdot \prod_{m=1}^{N-1}(1 - q^m)^{N-m}
  \]
  is a finite product: the \emph{Miura defect}. All genus-$1$
  arithmetic content in $\cW_N$ beyond Heisenberg is carried by
  this finite product. The prime-side Dirichlet analogue is
  $D_{\cW_N}^{\mathrm{prime}}(u)
  = -\zeta(u+1)\sum_{m=1}^{N-1}(N-m)\,m^{-u}$;
  the packet connection splits into $(N-1)$ Heisenberg copies plus
  a defect governing modes $m = 1, \ldots, N-1$.
\end{remark}

\begin{remark}[Shadow-moduli resolution]\label{rem:shadow-moduli}
  The shadow-moduli map $t\colon \mathfrak{H} \to \bC$ solves
  $\prod_j(1 - \lambda_j\,t(q))^{c_j}
  = \det(1 - K_q^{\conn})$
  with spectral atoms $\lambda_j$ and multiplicities $c_j$.
  The connected free energy factorises as
  $F_1^{\conn}(q) = G_\cA(t(q))$ with
  \[
    G_\cA(t)
    = \int \log(1 - \lambda t)\, d\rho_\cA(\lambda),
  \]
  determined by the spectral measure $\rho_\cA$ alone.
  The moments $\mu_r = \int \lambda^r\, d\rho_\cA$ satisfy the
  MC recursion degree by degree: the bar-complex differential
  equations control the spectral resolution of the genus-$1$
  amplitude.
\end{remark}

% ================================================================
% References
% ================================================================

\begin{thebibliography}{99}

\bibitem{BD04}
A.~Beilinson and V.~Drinfeld,
\textit{Chiral algebras},
Amer. Math. Soc. Colloq. Publ., vol.~51, AMS, 2004.

\bibitem{Bocherer85}
S.~B\"ocherer,
\"Uber die Fourier-Jacobi-Entwicklung Siegelscher Eisensteinreihen II,
\textit{Math. Z.} \textbf{189} (1985), 81--110.

\bibitem{CG17}
K.~Costello and O.~Gwilliam,
\textit{Factorization algebras in quantum field theory},
Vol.~1 and~2, Cambridge Univ. Press, 2017 and 2021.

\bibitem{Epstein03}
P.~Epstein,
Zur Theorie allgemeiner Zetafunctionen,
\textit{Math. Ann.} \textbf{56} (1903), 615--644.

\bibitem{Epstein07}
P.~Epstein,
Zur Theorie allgemeiner Zetafunctionen. II,
\textit{Math. Ann.} \textbf{63} (1907), 205--216.

\bibitem{FurusawaMorimoto}
M.~Furusawa and K.~Morimoto,
Refined global Gross--Prasad conjecture on special Bessel
periods and B\"ocherer's conjecture,
\textit{J. Eur. Math. Soc. (JEMS)} \textbf{23} (2021), 1295--1331.

\bibitem{IK04}
H.~Iwaniec and E.~Kowalski,
\textit{Analytic number theory},
American Mathematical Society Colloquium Publications, 53.
American Mathematical Society, Providence, RI, 2004.

\bibitem{FBZ04}
E.~Frenkel and D.~Ben-Zvi,
\textit{Vertex algebras and algebraic curves},
2nd ed., Math. Surv. Monogr., vol.~88, AMS, 2004.

\bibitem{Kac98}
V.~Kac,
\textit{Vertex algebras for beginners},
2nd ed., Univ. Lecture Ser., vol.~10, AMS, 1998.

\bibitem{Siegel51}
C.~L.~Siegel,
Indefinite quadratische Formen und Funktionentheorie. I, II,
\textit{Math. Ann.} \textbf{124} (1951/52), 17--54, 364--387.

\bibitem{TUY89}
A.~Tsuchiya, K.~Ueno, and Y.~Yamada,
Conformal field theory on universal family of stable curves with
gauge symmetries,
\textit{Adv. Stud. Pure Math.} \textbf{19} (1989), 459--566.

\bibitem{Verlinde88}
E.~Verlinde,
Fusion rules and modular transformations in 2D conformal field theory,
\textit{Nuclear Phys.~B} \textbf{300} (1988), 360--376.

\bibitem{Zagier08}
D.~Zagier,
Elliptic modular forms and their applications,
in \textit{The 1-2-3 of modular forms}, Springer, 2008, pp.~1--103.

\bibitem{Zhu96}
Y.~Zhu,
Modular invariance of characters of vertex operator algebras,
\textit{J. Amer. Math. Soc.} \textbf{9} (1996), 237--302.

\end{thebibliography}

\end{document}
